\documentclass[10pt]{article}
\usepackage{template}
\usepackage{amsmath,amssymb,microtype,tikz,xurl,needspace}
\usetikzlibrary{arrows.meta}
\setlength{\emergencystretch}{2em}
\SetKeys{date=2026-10-25,vol={2026-W43},title={每周重大新闻综述｜第43周},author={ChatGPT 编辑部},note={\small 观察期：2026 年 10 月 18—24 日。截稿：25 日上午。区分独立核实的事实、各方声明与编辑判断；数学预印本未经过本刊独立证明复核。}}
\newcommand{\weeklylead}[2]{\par\needspace{7\baselineskip}\bigskip{\Large\bfseries\color{newsblue}#1\par}{\footnotesize\color{newsgray}#2\par}\medskip\color{black}}
\newcommand{\weeklyfocus}[2]{\par\needspace{6\baselineskip}\bigskip{\large\bfseries\color{newsblue}#1\par}{\footnotesize\color{newsgray}#2\par}\smallskip\color{black}}
\newcommand{\weeklybrief}[2]{\par\needspace{5\baselineskip}\medskip{\normalsize\bfseries\color{newsblue}#1\par}{\footnotesize\color{newsgray}#2\par}\smallskip\color{black}}
\renewcommand{\refname}{参考资料}
\begin{document}
\begin{documents}
\raggedcolumns
\setlength{\parindent}{1.8em}
\setlength{\parskip}{3pt}
\weeklylead{霍尔木兹：战争阴影怎样变成油价与民生成本}{头条｜10月22—24日｜战争、航运与能源｜路透社、美联社、金融时报}
10月22日，伊朗伊斯兰革命卫队海军声称，波斯湾北部四艘涉嫌非法转运石油的船只在前一日遭打击起火，另有三艘船在霍尔木兹海峡以西触雷。路透社记录了这份声明，美联社则明确表示尚不能独立核实全部船只损毁情况。声明未提供足以确认全部船名、伤亡和损毁程度的材料。因此，本刊所能确认的是伊朗方面提出了七艘船只受损的主张，而不是七起事件的所有细节均已查明。\cite{hormuzReuters,hormuzAP}

海峡安全风险并非突然出现。《金融时报》根据航运信息指出，持续的袭击、骚扰和改道正在改变波斯湾出口石油的方式。军方的单日声明、海事机构逐步核实的事件记录、船东公告和保险公司的风险评估，分别属于不同口径，不能简单相加。即使两家媒体报道了相同数字，也可能来自同一个未经核实的消息源；交叉核查首先要判断信息是否真正独立。\cite{hormuzFT}

霍尔木兹海峡是波斯湾油气向外运输的重要通道。绕行陆上管道并非对所有产油国和油种都可行，管道容量、终点港口、储罐与炼厂条件也会形成约束。船舶受到威胁时，航运公司可能增加保费、船员风险补偿、等待时间和护航安排；部分货物改由船对船转运，又会占用额外船舶和作业时间。市场所谓的“运力紧张”，并不必然意味着全世界的船只数量突然减少，而可能是同一批船舶在相同时间内完成的有效运输工作减少。

\textbf{价格为什么会在实际短缺之前上涨？}原油期货价格包含对未来供应、库存与风险的预期。只要交易者认为下个月出现严重中断的概率增加，就可能要求风险溢价；但价格上涨本身并不能证明今天已经损失相同规模的产量。汇率、炼厂需求、产油国政策和金融头寸都会共同作用。分析时应把现货运输量、库存变化、远期价格和航运费用分开，而不是把一条油价曲线当作战争损失清单。

10月22日金融市场的反应尤其值得谨慎解读。美联社报道，当日美国纳斯达克指数下跌约2.6\%，标普500指数下跌约1.2\%，布伦特原油上涨约4\%，报每桶113.17美元。能源冲击可能推高通胀预期和融资成本，但人工智能企业的资本开支、债务负担与估值争议也在影响科技股。一个交易日的共同波动不构成单一因果关系的证明。\cite{marketAP}

伊朗方面表示愿在特定条件下恢复谈判，并把美国对伊朗港口的封锁与海峡重新开放联系起来。美方则维持区域军事部署和商船通行保护。联合国呼吁保护商业航运、重返外交途径。各方声明既表达政策立场，也服务于谈判筹码；“愿意谈”与“已经达成协议”之间仍隔着执行条款、核查机制和互信。\cite{iranTalks,hormuzAP}

另一项危险变化是威胁对象可能扩展到油气设施、海底电缆和数据中心。伊朗官员警告，若美国与以色列再次发动打击，可能扩大报复目标；这属于公开威胁，而非相关设施已被摧毁的事实。基础设施一旦受损，影响会沿着供电、通信、运输、食品和医疗系统传播，最终由大量不参与决策的平民承担成本。\cite{iranGuardian}

\begin{center}
\begin{tikzpicture}[>=Stealth,font=\scriptsize,every node/.style={align=center}]
\node[draw=blue!50,rounded corners,fill=blue!5,text width=2.5cm] (a) at (0,0) {袭击与航路风险\\（事件须独立核查）};
\node[draw=orange!50,rounded corners,fill=orange!6,text width=2.5cm] (b) at (3.2,0) {租船、保险、等待\\有效运力};
\node[draw=blue!50,rounded corners,fill=blue!5,text width=2.5cm] (c) at (0,-1.15) {库存与交付预期\\供给风险溢价};
\node[draw=orange!50,rounded corners,fill=orange!6,text width=2.5cm] (d) at (3.2,-1.15) {油价与终端成本\\企业、居民};
\draw[->,thick] (a)--(b);\draw[->,thick] (a)--(c);\draw[->,thick] (b)--(d);\draw[->,thick] (c)--(d);
\end{tikzpicture}

{\footnotesize\color{newsgray}图1｜编辑部原创机制示意，2026-10-25；非定量预测，不代表具体航运事件已核实。}
\end{center}

\textbf{观察与判断。}本周能够较有把握确认的是航运安全形势紧张、市场对供应中断的担忧增强，以及双方仍在公开提出外交条件。是否进入新一轮全面战争，不能凭威胁性讲话提前宣布。今后应持续核查船舶事故清单、海峡实际通行量、保险费率、正式谈判文件与军事行动。对这种高度政治化的战时新闻，承认尚不知道的事情，往往比给出一个过早的确定结论更有价值。

\weeklylead{俄乌再谈和平：停火、领土与安全保证的三重难题}{深度综述｜10月22—24日｜欧洲安全与人道局势}
10月23日，乌克兰总统泽连斯基表示，乌克兰、俄罗斯和美国的谈判人员拟于11月上半月在阿布扎比举行技术层级会谈，并希望随后能够安排领导人会面。俄方官员此前也表示原则上愿恢复三方接触。美联社与路透社对日程的描述略有差异：俄方提到11月10日以后，乌方强调上半月尽早举行。这些公开说法显示谈判可能重新启动，但不等于会议已经举行，更不意味着三方已经签署停火协议。\cite{ukraineReuters,ukraineAP}

真正难以解决的仍是领土和安全安排。俄罗斯要求乌克兰接受其对克里米亚及2022年宣布吞并的四个乌克兰州的控制，并提出中立和限制乌克兰军力等条件。乌克兰拒绝以领土让步作为和平的先决条件，要求可靠的长期安全保障。双方争议不只是地图上某条线的位置，还涉及战后军队部署、国际承认、居民权利和违反协议的责任。\cite{ukraineObstacles}

为什么连“先停火还是先谈总协议”都存在分歧？停止射击可以立即降低伤亡，为人道救援、交换战俘和修复设施创造空间；但若缺少监督与执行机制，双方又会担心对手利用间歇期重新部署兵力。乌克兰倾向先冻结战线再谈政治安排，俄罗斯则反对没有总体协议的停火。两种顺序会让不同一方承担不同风险，不能只靠一句“和平总是好的”消除这些现实约束。

安全保证也不是单一概念。它可能包括多年军援、训练、情报合作、防空体系、国际观察员、自动恢复制裁，以及在特定条件下的军事介入承诺。不同安排的可信度取决于预算、法律授权、指挥链、盟友政治意愿和执行成本。泽连斯基本周强调需要足以维持乌克兰军队十年的资金支持，说明基辅并不愿意把战后安全完全寄托于谈判对手的口头承诺。\cite{ukraineReuters}

美国同时是斡旋者、援助方和潜在安全担保者，这几个角色之间可能产生政策张力。华盛顿可能希望尽快降低战争与能源成本；乌克兰更重视协议能否阻止未来再次遭到攻击；欧洲国家则必须计算长期防务责任。谈判能否产生可信成果，不只取决于美俄领导人是否愿意会面，还取决于乌克兰是否实质参与，欧洲安全安排能否写成可执行条款。

外交接触与战争行动仍然并存。美联社报道，俄军在有关接触的消息前后继续对基辅及哈尔科夫附近实施袭击，造成平民伤亡。伤亡数字可能随搜救而修订，但对城市和能源设施的攻击并未因谈判意向而停止。将“同意讨论”包装为“战争结束”，会掩盖仍在承受风险的居民。\cite{ukraineAP}

本周还出现了局部人道安排。据《卫报》报道，双方同意自10月27日起在赫尔松地区奥列什基暂时停止军事行动，以便援助和撤离。白宫将斡旋归功于梅拉尼娅·特朗普，但相关人道协调此前已持续数周，功劳归属不宜只凭单方声明判定。由于实施日期尚未到来，本刊只能记录安排，不能宣称救援已经顺利完成。\cite{ukraineOleshky}

乌方同时对卫星通信速度限制和防空拦截弹交付提出担忧，美国方面此前有不同表述。通信对现代战场指挥、无人机和后勤十分重要，但公开材料不足以让本刊独立判断限制的技术成因，更不能据此断定任何一方蓄意制造某种战场结果。对于俄乌双方的战果声明，也应采用同样的独立核查标准。

\textbf{接下来该看什么？}真正有意义的进展应包括会谈正式举行、书面议程公布、停火监督方案和人道行动实际执行。长期和平还必须处理领土、战争责任、重建资金、人员返乡和安全保证。承认局部停火与全面政治和解的区别，并非降低对主权和国际法的重视，而是避免把复杂的制度问题误写为领导人之间的一次简单交易。

\weeklylead{从芯片到生物技术：美中投资安全规则的新边界}{深度综述｜10月22—24日｜美国财政部、中国商务部、路透社}
美国财政部10月22日发布对外投资安全计划的拟议规则，依据2025年通过的《综合对外投资国家安全法》（COINS Act）提出更新。现行计划已对美国主体投资某些中国相关半导体、人工智能和量子信息企业设置禁止或申报义务；新版拟议规则进一步明确高性能计算、超算、超音速系统等领域，并提议将生物技术纳入监管。财政部开放公众意见至12月7日。这是征求意见稿，而不是10月22日已经全面生效的新禁令。\cite{treasuryNPRM,treasuryProgram}

10月24日，中国商务部发言人表示坚决反对，认为美方扩大对华投资限制、泛化国家安全概念，并保留采取相应措施的权利。美方则声称政策目的是阻止美国资本与相关经营支持被用于其认定的军事、情报和监控技术。两份官方材料可以确认双方的政策立场，但不能由任何一方的政治修辞直接推出全部法律效果。\cite{mofcom,treasuryNPRM}

讨论这件事首先要分清三种工具：出口管制约束产品、软件和技术的跨境转移；对外投资审查约束本国投资者对特定境外企业的资本投入及相关经营参与；入境投资审查则评估外国资本收购本国资产可能带来的风险。三者可能服务于相似的安全目标，但对象、程序和例外规则并不相同。因此，说“美国企业以后不能与中国企业合作”，明显超出了本次拟议规则的范围。

投资也不只是汇款。创业资本可能同时提供董事会席位、管理经验、供应链渠道、人才网络和海外市场资源。美方关注的正是资本与非财务支持结合后对敏感技术发展的影响。反对者则指出，定义过宽会抬高合规成本、阻碍正常商业活动，甚至促使投资者减少对相关领域的尽职调查。争议的关键不在于安全与商业哪一个更值得重视，而是规则能否清楚识别需要防范的风险。

拟议规则区分需要通知财政部的交易和直接禁止的交易，讨论适用美国主体、受控境外实体、例外交易以及投资者的知识标准。路透社报道，公开交易证券、小额交易和某些药物许可安排可能存在例外，但最终适用范围必须以正式规则为准。对于企业，真正重要的是目标企业的具体活动、股权与控制关系、交易结构及尽职调查记录，而非只按公司名称或行业标签粗略判断。\cite{treasuryScope}

生物技术尤其需要精细定义。药物研发、临床诊断、基因编辑、病原体研究和生物制造具有不同的用途与风险。把它们统称为一个没有内部差别的“敏感行业”，可能误伤跨国药物研发与临床合作；但如果定义过窄，政策制定者担心的双重用途技术又可能不受约束。财政部征求意见的技术阈值、例外和执行成本，应当逐条阅读，而不是只看新闻标题里“扩大限制”四个字。

同一时期，中国也在修订技术出口目录。路透社10月18日报道，中国拟对部分太阳能制造技术实施更严格的出口管理。它与美国投资限制并不完全对称：一方主要监管美国资本对外流动，另一方涉及中国技术向外转移。把两项制度简单称为“互相报复”，还需要额外的时间和因果证据；更稳妥的说法是，两国都在强化对特定技术能力跨境流动的管理。\cite{chinaSolar}

政策的不确定性可能先于法规生效产生影响。基金可能延后投资决策，企业可能重新评估股权架构，科研团队可能修改合作条款。这些是合理的潜在传导机制，但本刊没有数据证明其实际规模。公开征求意见也意味着规则仍可修订，市场主体有机会就定义、例外和合规成本提出证据。

\textbf{本刊判断。}应继续追踪最终规则对技术类别的定义、申报与禁止的分界、例外交易、执法透明度以及中方是否采取具体措施。国家安全并不是可以自动取消所有经济成本分析的万能理由，开放市场也不是无视双重用途风险的充分理由。制度质量最终要由边界是否明确、执行是否可预测来检验。

\weeklyfocus{法国街头：示威人数的争议与公共服务的困局}{重点报道｜10月24日｜美联社、《卫报》}
10月24日，法国多地再次发生大规模示威，诉求包括学校资源、教师短缺、生活成本和社会不平等。美联社援引内政部估计，全国约41万人参加，其中巴黎约11万人；工会方面则估计接近80万人。《卫报》也记录了两种统计之间的显著差距。官方与组织者可能使用不同统计方法和范围，不能随意选择对某种立场有利的一项作为唯一事实。\cite{franceAP,franceGuardian}

学生最初关注校舍条件、教师和升学制度，工会与其他群体又将诉求扩展到工资、房租和公共服务。不同参与者未必拥有同一套政策方案，但可以共享对生活成本与资源分配的不满。这种运动不能只用“支持或反对总统”一个维度解释。

巴黎等城市出现冲突。报道提到有人投掷物品，警方使用催泪瓦斯、水炮等手段，拘捕人数接近一千。和平参与者、涉嫌暴力行为者和执法人员的行为必须分别评价。个别暴力事件不能抹去和平抗议的权利，合理社会诉求也不能为伤害他人提供豁免。具体责任应由独立调查而非政治立场预先裁决。\cite{franceAP}

财政空间与公共服务需求之间的张力是事件的重要背景。学校、交通、社会保障和债务利息都需要预算支持；当家庭感到购买力下降时，财政调整的成本分配就会成为政治问题。后续值得观察的是教育预算、谈判结果与执法调查，而不仅是单日示威人数。

\weeklyfocus{Linux 7.3：低显存 GPU、沙箱和可扩展调度}{工程观察｜10月18日正式发布｜kernel.org、LWN、KernelNewbies}
Linux 7.3于10月18日正式发布。LWN记录，该版本包含18,765个非合并提交，来自3,021名开发者，其中807人首次参与。与其只关注版本号，更值得了解它解决了哪些系统级问题。\cite{kernelOrg,kernelLWN}

对图形工作负载而言，低显存 GPU 的资源争用是一个有趣案例。多个程序竞争 VRAM 时，过去的内存保护策略可能使后来启动、但更需要资源的程序无法获得足够显存，只能退回使用系统内存映射区。新机制尝试逐出未受保护的缓冲区，为应获保护的分配提供空间。它针对的是显存紧张下的行为，不意味着所有游戏都会自动提高帧率。

内核还完成了 sched\_ext 的子调度器机制，允许不同 cgroup 由不同 BPF 调度策略管理。新加入的 \texttt{failfs} 伪文件系统与 \texttt{fchroot(2)} 系统调用提供另一种文件系统隔离思路：进程可以使绝对路径查找失败，只通过事先持有的目录文件描述符访问资源。这一机制需要正确的权限与描述符设计，并非调用一次就获得万能沙箱。\cite{kernelNewbies}

Btrfs、ext4、FUSE 和 SMB 服务端也获得性能或功能改进，虚拟化方面的 \texttt{userfaultfd} 支持增强了来宾内存工作集跟踪。实际部署仍需考虑发行版补丁、硬件驱动和回归测试；发行版的“Linux 7.3”并不必然与上游发布包完全相同。

\weeklyfocus{数学拾趣：随机整数分拆为什么偏爱某些余数？}{组合与概率｜10月22日 arXiv 预印本｜李继友、赵思成}
把正整数 $n$ 写成若干正整数之和，不考虑加数顺序，称为一个整数分拆。$5$ 有七种分拆：$5$、$4+1$、$3+2$、$3+1+1$、$2+2+1$、$2+1+1+1$、$1+1+1+1+1$。如果从七种中等概率选一种，奇数加数的个数严格多于偶数加数的情形有四种，所以概率是 $4/7$。问题看似简单，真正有趣的是当 $n$ 越来越大时，这个概率会趋向什么。

李继友、赵思成10月22日提交预印本，研究两条互不相交的正整数等差数列 $R=\{r+jm\}$ 与 $S=\{s+jm\}$ 中的加数出现次数。设 $a=r/m$、$b=s/m$，作者声称证明：当 $n$ 趋于无穷，$R$ 类加数数量超过 $S$ 类的概率趋向正则化不完全 Beta 函数 $I_{1/2}(a,b)$。对于奇偶分拆，取 $r=1,s=2,m=2$，便得到极限 $I_{1/2}(1/2,1)=1/\sqrt2$；它并不等于有限例子 $n=5$ 的 $4/7$。\cite{mathPartition}

论文还给出计数差经适当缩放后的分布极限，涉及两个独立 Gamma 随机变量的对数差。方法包括 Boltzmann 条件化、特征函数和一致 Fourier 估计，用以控制固定总和这一约束的影响。作者还称解决了此前提出的两个极限交换问题。

\textbf{证据等级：}这是尚待同行评议的 arXiv 预印本。本刊核查了题名、作者、提交日期和摘要，没有独立复核27页证明，也不把作者声称的定理等同于已获学界普遍确认的结论。

\weeklyfocus{125个超导量子比特，传送的是编码后的逻辑态}{量子物理｜《自然·物理学》10月22日论文}
《自然·物理学》10月22日发表一项表面码逻辑态传送实验。研究团队在125个超导量子比特处理器上，实现码距最高为7的表面码状态传送，并通过调整局域纠缠门观察传送成功与否的阈值行为。论文还研究了引入非经典资源后的阈值变化。\cite{quantumTeleport}

量子隐形传态不是把物质瞬间搬运，而是利用纠缠、测量结果和相应操作，在另一个物理载体上重建量子态。表面码则把逻辑信息分布在多个物理量子比特中，使部分局域错误可以被检测与纠正。两者结合的困难，是传送过程中如何保持编码结构，而不只是传送一个孤立的物理量子比特。

码距大致刻画保护逻辑信息所需的非平凡错误操作尺度；更高码距通常需要更多物理资源。实验中的阈值行为说明，逻辑传送的成败与纠缠强度之间存在可测的集体关系。但这仍是特定处理器与实验电路上的结果，不能直接推导出已经建成远距离、可扩展的容错量子网络。

论文已经过同行评议，并提供源数据；但可扩展性仍受噪声、门保真度、解码器与硬件互连限制。研究的意义是为未来分布式量子计算提供具体实验依据，而不是宣告相关工程难题已全部解决。

\weeklyfocus{数学拾趣：涂色游戏与稀疏矩阵核有什么关系？}{组合数学｜10月22日 arXiv 预印本｜韩宗建、梁亦然}
图的零强迫过程有一个直观规则：先将若干顶点涂黑，其余涂白；如果某个黑顶点恰好只有一个白邻居，就把那个邻居也涂黑。重复操作，若最终全图变黑，初始集合便是一种零强迫集合。这类涂色问题与稀疏矩阵的核和最小秩有深刻联系。

韩宗建、梁亦然10月22日提交预印本，研究给定图中某些顶点作为向量精确非零支集时，何时存在符合该图模式的对称矩阵，使该向量属于矩阵核。论文摘要指出，支集须满足一种称为 fort 的条件：支集外不能有恰好一个邻居位于支集内的顶点。\cite{mathZero}

直觉上，如果支集外某顶点只连接一个非零坐标，矩阵方程对应的一行就难以通过其他非零项抵消它。作者进一步研究复数域上系数空间的不可约性、余维与边界复形，并提出零强迫数的组合公式和算法。

这是新预印本，不是已经完成独立验证的理论体系。本刊仅介绍摘要明确陈述的主张与定义，未复核代数几何证明。

\weeklyfocus{新鸭绿江大桥：中朝经贸改善与外交象征}{东北亚外交｜10月24日｜路透社、美联社}
10月24日，朝鲜领导人金正恩经新鸭绿江大桥前往中国丹东，与中方高级官员王沪宁共同参加活动。路透社和美联社均报道了这次访问。大桥在多年延迟后于本月投入使用，有望改善重型车辆与集装箱运输条件；但桥梁开通不等于实际贸易量已增长某个百分比，通关制度、制裁与货物流量仍是关键。\cite{nkReuters,nkAP}

中国是朝鲜重要的经济和外交伙伴，而朝鲜近年来也与俄罗斯深化合作。边境纪念设施与工业场所的行程具有重申历史关系和经济联系的象征意义，却不能仅凭仪式推断未公布的经济或安全协议。

韩国希望中国帮助恢复半岛对话，美朝接触的前景则仍存在不确定性。外交活动释放的信号和正式谈判的可执行成果，需要分开观察。

\weeklybrief{化学｜网状配位策略与原子分散催化剂}{《自然·合成》10月22日}
Jung、Moon、Shim等发表论文，提出通过网状配位策略制备高负载、耐久的原子分散催化剂。初始结构的原子分散与反应条件下长期不团聚，是需要分别验证的两个问题。论文发表并不意味着已实现工业化量产。\cite{chemCatalyst}

\weeklybrief{金融｜AI投资逻辑将迎来财报检验}{路透社10月23日市场前瞻}
在能源价格上涨和科技股回撤之后，市场关注微软、Alphabet、亚马逊等企业财报，尤其是数据中心资本开支能否带来足够收入与现金流。单日股价下跌既不等于人工智能技术失败，也不证明此前所有估值合理；应区分技术能力、商业需求和融资条件。\cite{marketReuters}

\weeklybrief{科学史｜卡诺循环怎样从蒸汽机走向热力学}{《自然·物理学评论》10月22日}
Raffaele Pisano回顾卡诺理想循环从蒸汽机效率问题发展为热力学理论工具的历史。理想循环不是现实机器必须照搬的操作流程，而是用于理解可逆过程与效率上限的模型。工程中的“怎样提高效率”和理论中的“最多能够达到多高效率”并非同一个问题。\cite{carnot}

\par\bigskip{\large\bfseries\color{newsblue}编者校核说明}\par
本期观察窗口为 10 月 18—24 日。政治信源分别记录当事方立场与可独立核查的报道；战时事件与外交承诺不以单方声明视为已证实。数学两篇均为作者预印本，本刊未复核证明；同行评审论文也不意味着工程应用已成熟。图1为编辑部原创机制示意，不是现场摄影或定量预测。
\par\medskip{\large\bfseries\color{newsblue}参考资料与原始出处}\par
\begingroup\footnotesize\renewcommand{\refname}{}\bibliographystyle{unsrt}\bibliography{refs}\endgroup
\end{documents}
\end{document}
